\section{Contribution, engineering value and limits}\label{sec:r9-discussion}
The new evidence concerns three kinds of engineering decisions. A consumer-only requirement relaxation can inherit an anchor bound when the complete failure event is pointwise covered. A producer edit can retain that bound only for a proven invariant requested projection. A remaining request requires direct renewal or refusal. The domain quantities, units, source lineage and dependency witnesses make these decisions inspectable. This is stronger than checking whether a stored card's label matches a request, but it is also implementable by a complete procedural system.

The strongest procedural comparator demonstrates that fact. It is given all dependency knowledge, event comparisons, records and caches and reproduces the contract's authorization and numerical outputs. Accordingly, the independent value shown here belongs to \emph{explicit executable semantics}, rather than a graph data structure. The interface centralizes the interpretation and produces a checkable witness; no experiment establishes that engineers maintain it faster, that a graph database is required, or that generic programming cannot achieve it. Change reasoning and interval refinement already exist in engineering informatics. Our claim is their concrete integration with finite-sample event evidence and the measured inspection trade-offs, not invention of those general tools.

Direct renewal is also a serious competitor. It requires more event-record evaluations but supports additional requests and can cache subsequent queries. Semantic closure is useful when applicable anchors already exist and repeated relaxations matter; direct renewal is useful when sufficient applicable raw outputs exist or the anchor bound is conservative. An engineering implementation should permit both. Restricting the direct baseline's old-card access would weaken that comparison; the pre-confirmation parity amendment removes that restriction and preserves its prior supported services.

The ROI experiment demonstrates selective transfer rather than complete task recovery. A geometric optical quantity can remain invariant while target coverage and delivery cost change. The full reporting service remains uncertified at the primary allowance. Relaxing risk or deleting quality requirements defines a different service; the paper reports such alternatives explicitly instead of treating them as improvements to the original request. Fixed2's retained numerical trade-offs likewise prevent an unsupported claim of universal system-level superiority.

External validity remains limited. Synthetic palette landmarks are diagnostic targets, not natural surface defects. The virtual producer uses prescribed trajectories, declared perturbations and a frozen camera-work clock. Native CPU observations are quantized and depend on a desktop environment; they do not characterize an embedded camera's real-time operation. Real deployments can contain unmodeled lighting, texture, exposure-control, scheduling and mechanical couplings that invalidate a dependency projection. A metadata match cannot reveal such changes. New matched calibration, physical evidence and a suitable distribution-monitoring mechanism would be required before claiming physical reliability.

The three task families share one geometric source and executor. They establish a broader executable consumer interface than localization alone, but do not establish generalization to unrelated engineering domains. Deterministic malformed-source fixtures test specified refusal logic, not adversarial robustness in an open system. Complete-route performance remains historical descriptive evidence without mission-level task certification. These limitations are substantive: the study improves the evidence interpretation and cross-task validation, while the representation's broad originality and application reach remain questions for peer review.

\begin{figure}[!htbp]\centering\includegraphics[width=.94\linewidth]{r8/physical_platforms.pdf}
\caption{Four original real platform photographs preserved from the supplied research assets. They document engineering context and do not validate the newly simulated image, association or report events. Full original-resolution photographs and five prior simulation views remain in the historical release.}\label{fig:r9-platforms}\end{figure}
